\documentclass{article}

\usepackage{mathtools}
\usepackage{fullpage}
\usepackage{multirow}
\usepackage{float}

\title{Loci PETSc Interface}
\date{\nonumber}

\begin{document}

\maketitle

\section{Introduction}

This is re-implementation of the Loci PETSc interface for solution to
the linear system solvers on CPU and CUDA devices. The \texttt{FVMMod}
module in the Loci framework provides PETSc interface for linear
system solve. Following PETSc configurations are tested in this
interface.

\begin{table}[htbp]
\centering
\begin{tabular}{c|c|c|c}
\texttt{PetscInt} & \texttt{PetscScalar} & CPU & CUDA \\
\hline
32-bit & \texttt{double} & Yes & Yes \\
32-bit & \texttt{float} & Yes & Yes \\
\end{tabular}
\end{table}

\section{PETSc Configuration}
Some relevant configuration options for building PETSc are as follows.

\begin{enumerate}
\item \texttt{--with-mpi=1}: Enables MPI parallelization.
\item \texttt{--with-cuda=1}: Enables CUDA backend. It requires cuBLAS
  and cuSparse libraries. Other CUDA related options include
  \texttt{--with-cuda-dir}, \texttt{--with-cuda-pkg-config},
  \texttt{--with-cuda-include}, \texttt{--with-cuda-lib},
  \texttt{--with-cuda-arch}. See PETSc documentation for more details.
\item \texttt{--with-debugging=0} disables debugging version of the
  libraries. The default of this option is 1. Setting this option to 0
  is important for release version of PETSc build.
\item \texttt{--with-precision=single/double} sets the precision of
  the PETSc library. Default is \texttt{double}.
\item \texttt{--with-scalar-type=real/complex} Sets the numerical
  type, either \texttt{real} or \texttt{complex}. \texttt{real} is the
  default. This option must be \texttt{real}.
\end{enumerate}

\section{Loci PETSc Interface}
\texttt{FVMMod} provides parametric rules listed in
Table~\ref{tbl:Loci-PETSc-Rules} for interfacing with PETSc linear
system solvers. These rules depend on \texttt{X\_D}, \texttt{X\_L},
and \texttt{X\_U} facts that define the system matrix and
\texttt{X\_B} that defines the right hand side. The types of these
facts and the constraints on which these should be defined are listed
in Table~\ref{tbl:Loci-PETSc-Rule-Inputs}.
\begin{table}[H]
    \centering
    \begin{tabular}{c|c|c|c}
        \hline
        Parametric Rule & Type & Defined on & Use \\
        \hline
        \texttt{petscScalarSolve(X)} &  \texttt{store<double>} & \texttt{geom\_cells} & PETSc scalar solver \\
        \texttt{petscBlockedSolve(X)} &  \texttt{storeVec<double>} & \texttt{geom\_cells} & PETSc blocked solver \\
        \texttt{petscBlockedSSolve(X)} &  \texttt{storeVec<double>} & \texttt{geom\_cells} & PETSc blocked solver \\
    \end{tabular}
    \caption{Loci parametric rules for the PETSc interface.}
    \label{tbl:Loci-PETSc-Rules}
\end{table}
\begin{table}[H]
    \centering
    \begin{tabular}{c|c|c|c|c}
        Input & Defined on & \texttt{petscScalarSolve(X)} & \texttt{petscBlockedSolve(X)} & \texttt{petscBlockedSSolve(X)} \\
        \hline
        \texttt{X\_D} & \texttt{geom\_cells} & \texttt{store<double>} & \texttt{storeMat<double>} & \texttt{storeMat<float>} \\
        \texttt{X\_L} & \texttt{lower->cl} & \texttt{store<double>} & \texttt{storeMat<double>} & \texttt{storeMat<float>} \\
        \texttt{X\_U} & \texttt{upper->cr} & \texttt{store<double>} & \texttt{storeMat<double>} & \texttt{storeMat<float>} \\
        \texttt{X\_B} & \texttt{geom\_cells} & \texttt{store<double>} & \texttt{storeVec<double>} & \texttt{storeVec<float>} \\
    \end{tabular}
    \caption{Inputs to the Loci parametric rules for the PETSc interface.}
    \label{tbl:Loci-PETSc-Rule-Inputs}
\end{table}

For example, if a Loci solver needs to solve fixed point problem for a
system of PDEs, with variable \texttt{fluid}, and the solver
formulates and stores the resulting linear system in facts
\texttt{X\_D}, \texttt{X\_L}, \texttt{X\_U}, and \texttt{X\_B}, the
solver can invoke the PETSc linear system solver using following
parametric rule in a \texttt{.loci} file.
\begin{verbatim}
$type petscBlockedSSolve(fluid) storeVec<real>;
$rule pointwise(dfluid <- petscBlockedSSolve(fluid)), constraint(geom_cells) {
  $dfluid = $petscBlockedSSolve(fluid);
}
\end{verbatim}
The defaults for the Loci PETSc interface are listed in
Table\ref{tbl:Loci-PETSc-Defaults}.  With these defaults, the PETSc
solvers run on CPU, use sequential matrix and vector formats for
single process execution, and parallel matrix and vector formats for
parallel runs.
\begin{table}[H]
    \centering
    \begin{tabular}{c|c|c}
         & Scalar solve & Blocked solve \\
         \hline
        Matrix Type & \texttt{aij} & \texttt{baij} \\
        Vector Type & \texttt{standard} & \texttt{standard} \\
        Relative Tolerance & $1\times10^{-5}$ & $1\times10^{-5}$ \\
        Absolute Tolerance & $1\times10^{-18}$ & $1\times10^{-18}$ \\
        Maximum Iterations & 100 & 100 \\
        KSP Solver & PETSc Default & PETSc Default \\
        KSP Preconditioner & PETSc Default & PETSc Default \\
    \end{tabular}
    \caption{Defaults for Loci PETSc interface.}
    \label{tbl:Loci-PETSc-Defaults}
\end{table}

\section{Setting PETSc Options}
The PETSc interface provides \texttt{PETSCSolverControls} options list
to control behavior of PETSc. If the flow solver has multiple PDEs,
then each PDE linear system solver can be configured with different
PETSc options. For example, in turbulent flow simulation with SST
turbulence model, additional scalar transport equations are that of
turbulent kinetic energy ($k$) and turbulent dissipation rate
($\omega$). The \texttt{PETSCSolverControls} options list can be
specified as follows to set options for the individual variable.
\begin{verbatim}
PETSCSolverControls: <
  fluid=options(mat_type=baij, vec_type=standard),
  k=options(mat_type=aij, vec_type=standard, ksp_rtol=1e-3),
  w=options(mat_type=aij, vec_type=standard, ksp_rtol=1e-5)
>
\end{verbatim}
Note that, \texttt{k}, and \texttt{w} PETSc solvers use different
relative tolerance than the default as specified by option
\texttt{ksp\_rtol}. The names and values of the PETSc options are the
same as that in the PETSc documentation. This is demonstrated in the
following example.  The \texttt{fluid} PETSc solver preconditioner is
set to \texttt{sgs} and option \texttt{ksp\_converged\_reason} option
is used to print convergence reason.
\begin{verbatim}
PETSCSolverControls: <
  fluid=options(mat_type=baij, vec_type=standard, pc_type=sgs, ksp_converged_reason),
  k=options(mat_type=aij, vec_type=standard, ksp_rtol=1e-3),
  w=options(mat_type=aij, vec_type=standard, ksp_rtol=1e-5)
>
\end{verbatim}

\subsection{Offloading to CUDA Device}
To offload computation to CUDA devices, matrix and vector types need
to be set to the CUDA appropriate types. For example,
\begin{verbatim}
PETSCSolverControls: <
  fluid=options(mat_type=aijcusparse, vec_type=cuda),
  k=options(mat_type=aij, vec_type=standard, ksp_rtol=1e-3),
  w=options(mat_type=aij, vec_type=standard, ksp_rtol=1e-5)
>
\end{verbatim}
Note that, \texttt{k}, and \texttt{w} PETSc solvers run on CPU (as
indicated by their \texttt{mat\_type} and \texttt{vec\_type}) whereas
the fluid PDE linear system solver runs on CUDA device. Also note that
\texttt{mat\_type} is a scalar CUDA sparse type. The scalar type is
used in case of the fluid equations (which produced a blocked system)
because CUDA does not support blocked format for CUDA matrices.

\subsection{Specifying Global PETSc Options}
Global PETSc options are the options that are not specific to a
particular linear system. For example option \texttt{log\_view} prints
summary of log information, whereas \texttt{-log\_view
  :filename.txt:ascii\_flamegraph} saves the logging information in a
format suitable for visualizing as a Flame Graph. Option
\texttt{-log\_view\_memory} also displays memory usage in each
event. Option \texttt{-log\_view\_gpu\_time} displays time in each
event for GPU kernels, but slows down the computation.  These global
options are specified in \texttt{PETSCSolverControls} as follows.
\begin{verbatim}
PETSCSolverControls: <
  generic=options(log_view=":filename.txt:ascii_flamegraph", log_view_gpu_time)
>
\end{verbatim}

\section{Matrix and Vector Assembly Option}
PETSc provides two methods for assembling matrix and vectors.  In the
traditional method, a matrix is first preallocated using
\texttt{Mat***SetPreallocation()}. Then \texttt{MatSetValues()}
routines are invoked multiple times to specify element values,
followed by \texttt{MatAssemblyBegin()}, \texttt{MatBeginEnd()} to
initiate and finalize the communication needed to complete the
assembly. For vectors similar routines are used. PETSc introduced
coordinate assembly that reduces the cost of matrix and vector
assembly on GPUs. In this method, matrix is first preallocated using
\texttt{MatSetPreallocationCOO()} where the non-zero elements indices
are specified in coordinate format, followed by a single call to
\texttt{MatSetValuesCOO()} to set the element values. The order of the
element values must be the same as the order in which the element
coordinates are specified. By default, the Loci PETSc interface uses
the traditional method of matrix assembly. To enable the coordinate
assembly, use \texttt{petscCOOAssembly} in the vars file as follows.
\begin{verbatim}
petscCOOAssembly: 1
\end{verbatim}


\end{document}
